\documentclass[11pt,oneside]{article}

\pdfminorversion=7
\usepackage{QuizStyle}

\hypersetup{
    pdftitle={20260928 Quiz for MATH210 Applied Complex Variables},
    pdfauthor={Jungwoo Kim},
    pdfsubject={Quiz for MATH210 Applied Complex Variables},
    pdfcreator={LaTeX}
}

\begin{document}

\quiztitle{20260928 Quiz}

\begin{exercise}{26.4}
In each case, determine the singular points of the function and state why the function is analytic everywhere else:
\begin{exercisesubparts}
    \item \(\displaystyle f(z)=\frac{2z+1}{z(z^2+1)}\).
    \item \(\displaystyle f(z)=\frac{z^3+i}{z^2-3z+2}\).
    \item \(\displaystyle f(z)=\frac{z^2+1}{(z+2)(z^2+2z+2)}\).
\end{exercisesubparts}
\end{exercise}

\begin{solution}
Each function is rational, so by the differentiation rules it is analytic wherever its denominator is nonzero.
\begin{exercisesubparts}
    \item
    \[
        z(z^2+1)=z(z-i)(z+i),
    \]
    so the singular points are
    \[
        \boxed{z=0,\ \pm i}.
    \]

    \item
    \[
        z^2-3z+2=(z-1)(z-2),
    \]
    so the singular points are
    \[
        \boxed{z=1,\ 2}.
    \]

    \item Since
    \[
        z^2+2z+2=(z+1)^2+1=(z+1-i)(z+1+i),
    \]
    the singular points are
    \[
        \boxed{z=-2,\ -1\pm i}.
    \]
\end{exercisesubparts}
In each case the function is analytic at every other point.
\end{solution}

\Needspace{12\baselineskip}
\begin{exercise}{26.6}
Use results in Sec.~24 to verify that the function
\[
    g(z)=\ln r+i\theta
    \qquad(r>0,\ 0<\theta<2\pi)
\]
is analytic in the indicated domain of definition, with derivative \(g'(z)=1/z\). Then show that the composite function
\[
    G(z)=g(z^2+1)
\]
is analytic in the quadrant \(x>0\), \(y>0\), with derivative
\[
    G'(z)=\frac{2z}{z^2+1}.
\]
\begin{suggestion}
Observe that \(\ImPart(z^2+1)>0\) when \(x>0\), \(y>0\).
\end{suggestion}
\end{exercise}

\begin{solution}
For
\[
    g(z)=u(r,\theta)+iv(r,\theta),
    \qquad
    u=\ln r,\quad v=\theta,
\]
we have
\[
    r u_r=1=v_\theta,
    \qquad
    u_\theta=0=-r v_r.
\]
The relevant partial derivatives are continuous for \(r>0\), \(0<\theta<2\pi\). Hence the theorem in Sec.~24 gives
\[
    g'(z)=e^{-i\theta}(u_r+i v_r)
    =\frac{e^{-i\theta}}r
    =\boxed{\frac1z}.
\]
Now let \(z=x+iy\) with \(x>0\) and \(y>0\). Since
\[
    z^2+1=(x^2-y^2+1)+2xyi,
\]
we have
\[
    \ImPart(z^2+1)=2xy>0.
\]
Thus \(z^2+1\) lies in the domain of the chosen branch of \(g\). The inner function \(z^2+1\) is entire, so the composition is analytic in the quadrant. By the chain rule,
\[
    \boxed{G'(z)=g'(z^2+1)\cdot2z=\frac{2z}{z^2+1}}.
\]
\end{solution}

\begin{exercise}{26.7}
Let a function \(f\) be analytic everywhere in a domain \(D\). Prove that if \(f(z)\) is real-valued for all \(z\) in \(D\), then \(f(z)\) must be constant throughout \(D\).
\end{exercise}

\begin{solution}
Write
\[
    f(z)=u(x,y)+iv(x,y).
\]
Since \(f\) is real-valued throughout \(D\),
\[
    v(x,y)=0
\]
there, so \(v_x=v_y=0\). Because \(f\) is analytic, the Cauchy--Riemann equations give
\[
    u_x=v_y=0,
    \qquad
    u_y=-v_x=0.
\]
Consequently
\[
    f'(z)=u_x+i v_x=0
\]
throughout \(D\). By the theorem in Sec.~25, an analytic function whose derivative vanishes throughout a domain is constant. Hence \(f\) is constant on \(D\).
\end{solution}

\begin{exercise}{27.2}
Let the function \(f(z)=u(x,y)+iv(x,y)\) be analytic in a domain \(D\), and consider the families of level curves \(u(x,y)=c_1\) and \(v(x,y)=c_2\), where \(c_1\) and \(c_2\) are arbitrary real constants. Prove that these families are orthogonal. More precisely, show that if \(z_0=(x_0,y_0)\) is a point in \(D\) which is common to two particular curves \(u(x,y)=c_1\) and \(v(x,y)=c_2\) and if \(f'(z_0)\ne0\), then the lines tangent to those curves at \((x_0,y_0)\) are perpendicular.
\begin{suggestion}
Note how it follows from the pair of equations \(u(x,y)=c_1\) and \(v(x,y)=c_2\) that
\[
    u_x+u_y\frac{\dd y}{\dd x}=0
    \qquad\text{and}\qquad
    v_x+v_y\frac{\dd y}{\dd x}=0.
\]
\end{suggestion}
\end{exercise}

\begin{solution}
At \(z_0\), the Cauchy--Riemann equations give
\[
    u_x=v_y,
    \qquad
    u_y=-v_x.
\]
Also,
\[
    |f'(z_0)|^2=u_x^2+u_y^2,
\]
so \(f'(z_0)\ne0\) implies that \(u_x\) and \(u_y\) are not both zero.

Suppose first that \(u_xu_y\ne0\). Implicit differentiation gives the slopes
\[
    m_u=-\frac{u_x}{u_y},
    \qquad
    m_v=-\frac{v_x}{v_y}.
\]
Using the Cauchy--Riemann equations,
\[
    m_v=-\frac{-u_y}{u_x}=\frac{u_y}{u_x},
\]
and therefore
\[
    m_um_v=-1.
\]
Thus the tangent lines are perpendicular.

If \(u_y=0\), then \(u_x\ne0\). The level curve of \(u\) has a vertical tangent, while the Cauchy--Riemann equations give \(v_x=0\) and \(v_y=u_x\ne0\), so the level curve of \(v\) has a horizontal tangent. The case \(u_x=0\) is analogous. Hence the two families are orthogonal wherever \(f'(z_0)\ne0\).
\end{solution}

\begin{exercise}{29.1}
Use the theorem in Sec.~28 to show that if \(f(z)\) is analytic and not constant throughout a domain \(D\), then it cannot be constant throughout any neighborhood lying in \(D\).
\begin{suggestion}
Suppose that \(f(z)\) does have a constant value \(w_0\) throughout some neighborhood in \(D\).
\end{suggestion}
\end{exercise}

\begin{solution}
Suppose, to the contrary, that \(f(z)=w_0\) throughout some neighborhood contained in \(D\). Let
\[
    g(z)=w_0
\]
for all \(z\in D\). Then \(g\) is analytic in \(D\), and \(f\) and \(g\) agree throughout that neighborhood. By the uniqueness theorem in Sec.~28,
\[
    f(z)=g(z)=w_0
\]
throughout \(D\). This contradicts the assumption that \(f\) is not constant on \(D\). Therefore \(f\) cannot be constant throughout any neighborhood lying in \(D\).
\end{solution}

\begin{exercise}{29.5}
Show that if the condition that \(f(x)\) is real in the reflection principle (Sec.~29) is replaced by the condition that \(f(x)\) is pure imaginary, then equation (1) in the statement of the principle is changed to
\[
    \overline{f(\overline z)}=-f(z).
\]
\end{exercise}

\begin{solution}
Define
\[
    h(z)=-if(z).
\]
If \(f(x)\) is pure imaginary on the segment of the real axis, then \(h(x)\) is real there. Since \(h\) is analytic in the same domain, the reflection principle applied to \(h\) gives
\[
    \overline{h(\overline z)}=h(z).
\]
Substituting \(h=-if\),
\[
    i\,\overline{f(\overline z)}=-if(z),
\]
so
\[
    \boxed{\overline{f(\overline z)}=-f(z)}.
\]
Conversely, if this identity holds, then for real \(x\),
\[
    \overline{f(x)}=-f(x),
\]
which means that the real part of \(f(x)\) is zero. Hence \(f(x)\) is pure imaginary, as required.
\end{solution}

\begin{exercise}{30.5}
Write \(|\exp(2z+i)|\) and \(|\exp(iz^2)|\) in terms of \(x\) and \(y\). Then show that
\[
    |\exp(2z+i)+\exp(iz^2)|\le e^{2x}+e^{-2xy}.
\]
\end{exercise}

\begin{solution}
For any complex number \(w\),
\[
    |e^w|=e^{\RePart w}.
\]
Since \(z=x+iy\),
\[
    \RePart(2z+i)=2x,
\]
so
\[
    |\exp(2z+i)|=e^{2x}.
\]
Also,
\[
    z^2=x^2-y^2+2xyi,
\]
and hence
\[
    iz^2=-2xy+i(x^2-y^2).
\]
Therefore
\[
    |\exp(iz^2)|=e^{-2xy}.
\]
The triangle inequality now gives
\[
    \boxed{|\exp(2z+i)+\exp(iz^2)|\le e^{2x}+e^{-2xy}}.
\]
\end{solution}

\begin{exercise}{30.13}
Let the function \(f(z)=u(x,y)+iv(x,y)\) be analytic in some domain \(D\). State why the functions
\[
    U(x,y)=e^{u(x,y)}\cos v(x,y),
    \qquad
    V(x,y)=e^{u(x,y)}\sin v(x,y)
\]
are harmonic in \(D\).
\end{exercise}

\begin{solution}
Since \(f\) is analytic in \(D\) and the exponential function is entire, the composition \(e^{f(z)}\) is analytic in \(D\). Moreover,
\[
\begin{aligned}
    e^{f(z)}
    &=e^{u+iv}\\
    &=e^u(\cos v+i\sin v)\\
    &=U+iV.
\end{aligned}
\]
By the theorem in Sec.~27, the real and imaginary components of an analytic function are harmonic. Therefore \(U\) and \(V\) are harmonic throughout \(D\).
\end{solution}

\begin{exercise}{33.2}
Show that
\begin{exercisesubparts}
    \item \(\log e=1+2n\pi i\), \(n=0,\pm1,\pm2,\ldots\);
    \item \(\displaystyle \log i=\left(2n+\frac12\right)\pi i\), \(n=0,\pm1,\pm2,\ldots\);
    \item \(\displaystyle \log(-1+\sqrt3\,i)=\ln2+2\left(n+\frac13\right)\pi i\), \(n=0,\pm1,\pm2,\ldots\).
\end{exercisesubparts}
\end{exercise}

\begin{solution}
For \(z\ne0\),
\[
    \log z=\ln|z|+i(\Arg z+2n\pi),
    \qquad n\in\Z.
\]
\begin{exercisesubparts}
    \item Since \(|e|=e\) and \(\Arg e=0\),
    \[
        \boxed{\log e=1+2n\pi i}.
    \]

    \item Since \(|i|=1\) and \(\Arg i=\pi/2\),
    \[
        \boxed{\log i=i\left(\frac\pi2+2n\pi\right)
        =\left(2n+\frac12\right)\pi i}.
    \]

    \item The number \(-1+\sqrt3\,i\) has modulus \(2\) and principal argument \(2\pi/3\). Thus
    \[
    \begin{aligned}
        \log(-1+\sqrt3\,i)
        &=\ln2+i\left(\frac{2\pi}{3}+2n\pi\right)\\
        &=\boxed{\ln2+2\left(n+\frac13\right)\pi i}.
    \end{aligned}
    \]
\end{exercisesubparts}
\end{solution}

\begin{exercise}{33.4}
Show that
\[
    \log(i^2)\ne2\log i
\]
when the branch
\[
    \log z=\ln r+i\theta
    \qquad
    \left(r>0,\ \frac{3\pi}{4}<\theta<\frac{11\pi}{4}\right)
\]
is used. (Compare this with the example in Sec.~33.)
\end{exercise}

\begin{solution}
Let \(g\) denote the single-valued logarithmic branch specified in the exercise, written as \(\log\) in the textbook. For \(i\), the arguments are \(\pi/2+2k\pi\). The unique one in the branch interval
\[
    \frac{3\pi}{4}<\theta<\frac{11\pi}{4}
\]
is \(5\pi/2\). Hence
\[
    g(i)=\frac{5\pi}{2}i,
    \qquad
    2g(i)=5\pi i.
\]
On the other hand, \(i^2=-1\), and the unique argument of \(-1\) in the same interval is \(\pi\). Thus
\[
    g(i^2)=g(-1)=\pi i.
\]
Therefore
\[
    \boxed{g(i^2)=\pi i\ne5\pi i=2g(i)}.
\]
\end{solution}

\begin{exercise}{33.11}
Show in two ways that the function
\[
    \ln(x^2+y^2)
\]
is harmonic in every domain that does not contain the origin.
\end{exercise}

\begin{solution}
Let
\[
    H(x,y)=\ln(x^2+y^2).
\]
First, direct differentiation gives
\[
    H_x=\frac{2x}{x^2+y^2},
    \qquad
    H_y=\frac{2y}{x^2+y^2},
\]
and, for \((x,y)\ne(0,0)\),
\[
    H_{xx}=\frac{2(y^2-x^2)}{(x^2+y^2)^2},
    \qquad
    H_{yy}=\frac{2(x^2-y^2)}{(x^2+y^2)^2}.
\]
Hence
\[
    H_{xx}+H_{yy}=0.
\]
All first- and second-order partial derivatives are continuous away from the origin, so \(H\) is harmonic there.

For a second argument, fix any \(z_0\ne0\). A sufficiently small disk about \(z_0\) avoids the origin and admits an analytic branch \(g\) of the logarithm. For a continuous choice \(\theta(z)\) of argument on that disk,
\[
    g(z)=\ln|z|+i\theta(z),
    \qquad
    \RePart g(z)=\ln|z|=\frac12\ln(x^2+y^2).
\]
By the theorem in Sec.~27, the real part of an analytic function is harmonic. Consequently \(H=2\RePart g\) is harmonic near \(z_0\). Since \(z_0\) was arbitrary, \(H\) is harmonic throughout every domain that excludes the origin; a global logarithm branch on that domain is not required.
\end{solution}

\begin{exercise}{34.1}
Show that for any two nonzero complex numbers \(z_1\) and \(z_2\),
\[
    \Log(z_1z_2)=\Log z_1+\Log z_2+2N\pi i,
\]
where \(N\) has one of the values \(0,\pm1\). (Compare with Example 2 in Sec.~34.)
\end{exercise}

\begin{solution}
Let
\[
    \theta_1=\Arg z_1,
    \qquad
    \theta_2=\Arg z_2,
    \qquad
    S=\theta_1+\theta_2.
\]
Since \(-\pi<\theta_j\le\pi\),
\[
    -2\pi<S\le2\pi.
\]
The principal argument of the product is obtained by bringing \(S\) into \((-\pi,\pi]\):
\[
    \Arg(z_1z_2)=S+2N\pi,
\]
where
\[
    N=
    \begin{cases}
        -1, & S>\pi,\\
        0, & -\pi<S\le\pi,\\
        1, & S\le-\pi.
    \end{cases}
\]
Thus \(N\in\{0,\pm1\}\). Since \(|z_1z_2|=|z_1||z_2|\),
\[
\begin{aligned}
    \Log(z_1z_2)
    &=\ln|z_1z_2|+i\Arg(z_1z_2)\\
    &=\ln|z_1|+\ln|z_2|+i(\theta_1+\theta_2+2N\pi)\\
    &=\boxed{\Log z_1+\Log z_2+2N\pi i}.
\end{aligned}
\]
\end{solution}

\begin{exercise}{34.2}
Verify expression (4), Sec.~34, for \(\log(z_1/z_2)\) by
\begin{exercisesubparts}
    \item using the fact that \(\arg(z_1/z_2)=\arg z_1-\arg z_2\) (Sec.~9);
    \item showing that \(\log(1/z)=-\log z\), \(z\ne0\), in the sense that \(\log(1/z)\) and \(-\log z\) have the same set of values, and then referring to expression (1), Sec.~34, for \(\log(z_1z_2)\).
\end{exercisesubparts}
\end{exercise}

\begin{solution}
The identities here are identities of multiple-valued logarithms, so equality means equality of the corresponding sets of values, with compatible values selected.
\begin{exercisesubparts}
    \item Choose arguments \(\theta_1\in\arg z_1\) and \(\theta_2\in\arg z_2\). The argument identity in Sec.~9 states that all arguments of \(z_1/z_2\) have the form \(\theta_1-\theta_2+2k\pi\), where \(k\in\Z\). Therefore all logarithmic values of the quotient are
    \[
        \ln\frac{|z_1|}{|z_2|}
        +i(\theta_1-\theta_2+2k\pi),
        \qquad k\in\Z.
    \]
    These are precisely the differences of the values
    \(\ln|z_1|+i(\theta_1+2m\pi)\) and
    \(\ln|z_2|+i(\theta_2+2n\pi)\), because \(m-n\) ranges over \(\Z\). Hence, as sets of values,
    \[
        \boxed{\log\left(\frac{z_1}{z_2}\right)=\log z_1-\log z_2}.
    \]

    \item Write \(z=re^{i\theta}\). The values of \(\log(1/z)\) are
    \[
        -\ln r+i(-\theta+2k\pi),
        \qquad k\in\Z.
    \]
    The values of \(-\log z\) are
    \[
        -\ln r-i(\theta+2n\pi)
        =-\ln r+i(-\theta-2n\pi),
        \qquad n\in\Z.
    \]
    These are the same set of values, so
    \[
        \log(1/z)=-\log z.
    \]
    Since
    \[
        \frac{z_1}{z_2}=z_1\frac1{z_2},
    \]
    expression (1), Sec.~34, gives
    \[
        \log\left(\frac{z_1}{z_2}\right)
        =\log z_1+\log(1/z_2)
        =\boxed{\log z_1-\log z_2}.
    \]
\end{exercisesubparts}
\end{solution}

\Needspace{23\baselineskip}
\begin{exercise}{34.5}
Let \(z\) denote any nonzero complex number, written
\[
    z=re^{i\Theta}
    \qquad(-\pi<\Theta\le\pi),
\]
and let \(n\) denote any fixed positive integer. Show that all of the values of \(\log(z^{1/n})\) are given by
\[
    \log(z^{1/n})
    =\frac1n\ln r
    +i\frac{\Theta+2(pn+k)\pi}{n},
\]
where \(p=0,\pm1,\pm2,\ldots\) and \(k=0,1,2,\ldots,n-1\). Then, after writing
\[
    \frac1n\log z
    =\frac1n\ln r+i\frac{\Theta+2q\pi}{n},
    \qquad q=0,\pm1,\pm2,\ldots,
\]
show that the set of values of \(\log(z^{1/n})\) is the same as the set of values of \((1/n)\log z\). Thus show that
\[
    \log(z^{1/n})=\frac1n\log z,
\]
where, corresponding to a value of \(\log(z^{1/n})\) taken on the left, the appropriate value of \(\log z\) is to be selected on the right, and conversely. [The result in Exercise 5, Sec.~33, is a special case of this one.]
\begin{suggestion}
Use the fact that the remainder upon dividing an integer by a positive integer \(n\) is always an integer between \(0\) and \(n-1\), inclusive; that is, any integer \(q\) can be written \(q=pn+k\), where \(p\in\Z\) and \(k\in\{0,1,\ldots,n-1\}\).
\end{suggestion}
\end{exercise}

\begin{solution}
The \(n\) distinct values of \(z^{1/n}\) are
\[
    w_k=r^{1/n}\exp\left(i\frac{\Theta+2k\pi}{n}\right),
    \qquad k=0,1,\ldots,n-1.
\]
For each fixed \(k\), all values of \(\log w_k\) are
\[
\begin{aligned}
    \log w_k
    &=\frac1n\ln r
      +i\left(\frac{\Theta+2k\pi}{n}+2p\pi\right)\\
    &=\frac1n\ln r
      +i\frac{\Theta+2(pn+k)\pi}{n},
      \qquad p\in\Z.
\end{aligned}
\]
Thus these formulas give all values of \(\log(z^{1/n})\).

On the other hand, the values of \((1/n)\log z\) are
\[
    \frac1n\ln r+i\frac{\Theta+2q\pi}{n},
    \qquad q\in\Z.
\]
By the division algorithm, every integer \(q\) can be written uniquely as
\[
    q=pn+k,
    \qquad p\in\Z,
    \qquad 0\le k\le n-1.
\]
Hence the two displayed families of values are exactly the same. Therefore, in the multiple-valued sense described in Sec.~34,
\[
    \boxed{\log(z^{1/n})=\frac1n\log z}.
\]
\end{solution}

\end{document}
